\subsection{A torsion bound for mixed full-support shallow lattices}

\begin{lemma}[Proved: commuting translations at an arbitrary cubic background]
\label{C:shallow-commuting-translations}
Let $A=\mathbf Z_3[\zeta]$, $\pi=\zeta-1$, and
$k=A\ltimes J$, where $J$ is an arbitrary $\mathbf F_3$-vector
space.  Fix a finite free constant cubic lattice $L$ with
matrix $T_0$, and write
\[
 d(U)=[(L\otimes_Ak,T_0+U)]-[(L\otimes_Ak,T_0)]
\]
for any $J$-valued tangent perturbation defining a cubic action.
Put $z=[k_\zeta]$.  For a fixed such $U$, define
\[
 f_U(H)=d(U+H)-d(U)
 \quad\text{for }H\in M_{\operatorname{rank}L}(J)
          \text{ commuting with }\overline T_0.
\]
Then $f_U$ is additive, $z f_U(H)=f_U(H)$, and
$3f_U(H)=0$.
These conclusions do not assert that $f_U(H)=d(H)$.
\end{lemma}
\begin{proof}
Every commuting $H$ is itself tangent, since the derivative
of the cubic relation is $3\overline T_0^{2}H=0$.
For tangent $U,V$ consider
\[
 E(U,V)=\begin{pmatrix}T_0+U&I\\0&\zeta T_0+V\end{pmatrix}.
\]
Its diagonal cubes are identities.  The constant upper-right
block of its cube is zero.  Since products of two $J$-valued
entries vanish and $\zeta J=J$, its linear part is
\[
 \overline T_0U+2U\overline T_0
          +2\overline T_0V+V\overline T_0
       =[\overline T_0,U-V].
\]
Thus $E(U,V)$ is a cubic action whenever $U-V$ commutes
with $\overline T_0$.
For commuting $H$, conjugation by
$\left(\begin{smallmatrix}I&0\\H&I\end{smallmatrix}\right)$
then gives exactly
\[
 \begin{pmatrix}T_0+U+H&I\\0&\zeta T_0+V-H\end{pmatrix}.
\]
The lower-left block is $[\overline T_0,H]=0$.
Both triangular matrices yield admissible free-block exact
sequences, so
\[
 d(U)+z d(V)=d(U+H)+z d(V-H).                    \tag{C.1.MT1}
\]
Take $V=U+K$ with $K$ commuting.  Subtracting
$(1+z)d(U)$ gives
\[
 z f_U(K)=f_U(H)+z f_U(K-H).
\]
Putting $K=H$ proves $zf_U(H)=f_U(H)$, after which the
same equality proves additivity.  Its domain is an additive
$\mathbf F_3$-vector space, hence additivity gives
$3f_U(H)=f_U(3H)=0$.
No additivity in the noncommuting background $U$ was used.
\end{proof}

\begin{theorem}[Proved: three-torsion for arbitrary full-support cubic mixtures]
\label{C:shallow-full-support-three-torsion}
Keep the preceding notation.  Suppose the constant cubic
lattice $L$ is any direct sum of copies of $L_{012}$, $Q$,
and $A[C_3]$ from
Proposition~\ref{C:cyclotomic_C3_lattices}.
Every shallow deformation of $L$ then satisfies
\[
                         3d(U)=0
                         \quad\text{in }R_k(C_3).       \tag{C.1.MT2}
\]
In particular let $M$ be a free $k[S_3]$-representation whose
constant reduction is an arbitrary mixture of the four
signed $L_{012}$ and $Q$ types and the two signed regular
types of Proposition~\ref{C:cyclotomic_S3_lattices}.
Then
\[
 3\bigl([M]-[M/JM\otimes_Ak]\bigr)=0
                         \quad\text{in }R_k(S_3).       \tag{C.1.MT3}
\]
The theorem gives a torsion bound.  It does not assert
vanishing for the mixed types.
\end{theorem}
\begin{proof}
Constant regular cubic summands split off from every shallow
deformation by the equivariant coinduction argument in
Theorem~\ref{C:S3_isotypic_shallow_cancellation}.
They contribute zero relative class.
Thus it suffices to take
$L=L_{012}^{r}\oplus Q^{s}$.
Set $a=r+s$, and order the residual basis by its chain
vectors $E,F$ and remaining vectors $Y$.
In the integral bases already used in that theorem,
\[
 N=\overline T_0-I=
       \begin{pmatrix}0&I_a&0\\0&0&0\\0&0&0\end{pmatrix}.
\]
Lemma~\ref{C:shallow_tangent_quotient} allows us, by an
infinitesimal change of basis, to assume
\[
 U=\begin{pmatrix}0&0&0\\0&S&A\\B&0&C\end{pmatrix}.
                                                               \tag{C.1.MT4}
\]
All four lettered blocks have entries in $J$.

\emph{The constant character rotation.}
In normalized character coordinates let $P$ cyclically
permute the three coordinates by
$(a_0,a_1,a_2)\mapsto(a_1,a_2,a_0)$.
It preserves both the common-residue-line lattice $L_{012}$
and the residual-sum-zero-plane lattice $Q$.
For $D=\operatorname{diag}(1,\zeta,\zeta^2)$ one has
$P^{-1}(\zeta D)P=D$.
In the integral bases
\[
 \begin{array}{ll}
 L_{012}:&\pi(e_1-e_2),\ e_0+e_1+e_2,\ \pi e_0,\\
 Q:&\pi(e_1+e_2),\ e_1-e_2,\ e_0-(e_1+e_2)/2,
 \end{array}
\]
its respective residual matrices are $I+E_{13}$ and
$I+E_{32}$.
These formulas follow by applying the displayed coordinate
permutation to each of the three basis vectors; in particular
no abstract isomorphism between the two different lattice
types is being assumed.
For all copies together, write
\[
 \overline P=
 \begin{pmatrix}I&0&C_0\\0&I&0\\0&D_0&I\end{pmatrix},
 \qquad
 C_0=\operatorname{diag}(I_r,0_s),\quad
 D_0=\operatorname{diag}(0_r,I_s).
                                                               \tag{C.1.MT5}
\]
Here $C_0D_0=D_0C_0=0$.
Conjugating the $\zeta$-twist of the deformed action by
this integral $P$ gives $T_0+\overline P^{-1}U\overline P$.
Consequently
\[
                 z d(U)=d(\overline P^{-1}U\overline P).
                                                               \tag{C.1.MT6}
\]

\emph{Rotation changes only commuting tangent coordinates.}
Multiply (C.1.MT4) by (C.1.MT5), and take the four unique
tangent quotient coordinates of
Lemma~\ref{C:shallow_tangent_quotient}.
They are
\[
 \begin{split}
 S'&=S+AD_0-C_0B, & A'&=A,\\
 B'&=B,          & C'&=C+BC_0-D_0A.
 \end{split}                                                   \tag{C.1.MT7}
\]
Thus the difference between its gauge class and that of $U$
is represented by a commuting matrix: if
$\Delta S=AD_0-C_0B$ and $\Delta C=BC_0-D_0A$, choose
\[
 H=\operatorname{diag}(\Delta S/2,\Delta S/2,\Delta C).
\]
This commutes with $N$, and its quotient coordinates are
$(\Delta S,0,0,\Delta C)$.
By the uniqueness in the tangent quotient lemma, an actual
infinitesimal conjugacy identifies the perturbations
$\overline P^{-1}U\overline P$ and $U+H$.
Hence (C.1.MT6) gives
\[
             h:=(z-1)d(U)=f_U(H),\qquad zh=h,
             \qquad 3h=0,                                \tag{C.1.MT8}
\]
by the preceding lemma.
This step uses a commuting translation at the fixed
background, not a splitting of $d(U)$ into its four blocks.

The sum $1+z+z^2$ is the class of the constant regular
cubic lattice.  Indeed its monic character-factor filtration
over $A$ has free successive quotients, and therefore
remains admissible after base change to $k$.
Tensoring a free cubic representation with $k[C_3]$
untwists its action to the regular action on each underlying
basis vector.  Since the deformation and its constant
reduction have the same rank, it follows that
\[
                       (1+z+z^2)d(U)=0.
\]
From (C.1.MT8), $zd(U)=d(U)+h$ and
$z^2d(U)=d(U)+2h$.
The last displayed identity is therefore
$0=3d(U)+3h=3d(U)$, proving (C.1.MT2).

For the $S_3$ assertion apply this to the cubic restriction.
The relative class is generically zero, because
$J\otimes\mathbf Q=0$, and
$K_0(k)=\mathbf Z\to K_0(k\otimes\mathbf Q)=\mathbf Z$
is injective because both rings are local and nonzero.
The cyclic detection lemma
(Lemma~\ref{C:S3_cyclic_detection}) states that restriction
is injective on this generic kernel.
It therefore sends the equality just proved back to
(C.1.MT3).  No conclusion about vanishing of a possible
three-torsion mixed class follows from this argument.
\end{proof}
